% acknowledgements: written in Italian in both versions; generated from ../../ringraziamenti.txt
% one paragraph per person or group, separated by vertical space, recipient in bold
{\setlength{\parindent}{0pt}\setlength{\parskip}{0.9\baselineskip}

Il primo grazie va a \textbf{Torino}. Non è scontato che una città diventi casa, eppure Torino ci è riuscita quasi senza che me ne accorgessi: mi ha accolto con discrezione, senza clamore, e fin dai primi giorni mi ha fatto sentire nel posto giusto. Le sue strade con i suoi portici eleganti hanno fatto da sfondo a tutto quello che ho vissuto in questi anni, alle giornate di studio come alle serate che non dimenticherò. Ma una città, in fondo, è fatta soprattutto delle persone che la abitano. Qui ho incontrato chi, giorno dopo giorno, è diventato la mia seconda famiglia: persone con cui ho condiviso tavole, fatiche, risate e silenzi, che ci sono state nei momenti più leggeri e soprattutto in quelli più difficili, quando casa era lontana e loro erano lì. A loro devo molto più di quanto queste righe possano dire.

\textbf{Ai miei genitori}, i miei due fari, che anche da lontano non hanno mai smesso di indicarmi la rotta. Tutto quello che sono oggi lo devo a voi: in questi anni avete saputo guidarmi senza mai impormi una strada, rispettare le mie scelte anche quando vi costavano qualcosa e starmi accanto nei momenti in cui da solo non ce l'avrei fatta. Grazie per avermi cresciuto nel modo migliore in cui si potesse fare, per non avermi mai fatto mancare nulla e, soprattutto, per i valori che mi avete trasmesso e che ogni giorno, spesso senza che me ne accorga, orientano le mie decisioni. Grazie per avermi lasciato partire per Torino a inseguire i miei sogni, pur sapendo che quella partenza avrebbe messo tanti chilometri tra noi. Quei chilometri non ho mai imparato a sentirli: siete sempre stati al mio fianco, a ogni passo. Vi amo.

\textbf{A Benny}, da sempre la luce della nostra famiglia. Grazie per tutti i momenti passati insieme, per le risate, per quella spensieratezza che riesci a trasmettere a chiunque ti stia vicino e per la spontaneità che ti rende inconfondibile. Crescere insieme è stata una fortuna. Da poco anche tu hai iniziato la tua esperienza da fuorisede: è l'inizio di una nuova avventura, ed è giusto che sia tutta tua. Ti auguro con tutto il cuore di trovare la tua strada, qualunque essa sia. Ci saranno giornate entusiasmanti e altre in cui casa ti sembrerà lontanissima; in quelle, ricordati che io ci sono e ci sarò sempre, a qualunque distanza, ogni volta che avrai bisogno di me. E non permettere mai a nessuno di cambiare quello che sei.

\textbf{Alla mia famiglia tutta}: agli zii, ai nonni e ai cugini. Ognuno di voi, a modo suo, è per me un esempio e un punto di riferimento, e in ogni ritorno a casa ho ritrovato il calore che solo una famiglia sa dare. Grazie per il vostro sostegno, costante e silenzioso, che non mi è mai mancato.

\textbf{Ad Angelo}, il fratello inaspettato che Torino mi ha regalato. In brevissimo tempo sei diventato uno dei miei punti di riferimento laddove per me era tutto nuovo. Da amico di amici a compagno, in un attimo. Anche nelle situazioni difficili sei sempre riuscito a farmi trovare la giusta direzione. Insieme abbiamo già condiviso tantissimi momenti indimenticabili e sono certo che il tempo ce ne regalerà ancora molti.

\textbf{A Mary}. Con te tutto è sembrato facile da subito. Grazie per gli inciuci, i giovedì sera, gli sguardi d'intesa nelle situazioni divertenti, gli audio inutilmente lunghi, le mattine al mercato e tutto il resto. Grazie per esser sempre stata al mio fianco anche nei momenti più duri. Con te sono sempre riuscito a mettere in fila i pensieri, anche quando questi erano negativi. Spero che il nostro rapporto possa sempre continuare a evolvere come ha fatto finora.

\textbf{A Pietro}. Il ricordo di te che mi vieni incontro sulle scale alle aule R è ormai indelebile. Quello è stato per me il vero inizio di questa avventura, e per questo non ti ringrazierò mai abbastanza. Non ti ho mai nascosto la stima che provo nei tuoi confronti e sono sicuro che riuscirai a ottenere i risultati che ti prefiggerai. Ti auguro tutto il meglio che la vita potrà regalarti.

\textbf{A Lollo}. In questo momento mi sto chiedendo per cosa esattamente io debba ringraziarti, e in piena onestà\dots{} boh, non lo so. Ma la verità è che è proprio questo a renderti unico. La tua presenza riesce a essere sempre rumorosa, anche se stai seduto nello stesso punto per tutta la serata. Ti si ama così come sei, e questo non te lo toglierà nessuno. Ti auguro il meglio (\emph{skill issue}).

\textbf{A Mauro}. Anche se siamo molto diversi, in fondo in fondo sappiamo di essere anche molto simili, e questo ci ha portati a costruire un rapporto speciale a modo suo. Spesso sei l'unico che capisce le mie citazioni di AGG, quindi ormai sono diventate il nostro modo di comunicare, anche se all'esterno ci fanno apparire come due idioti (il che, comunque, non è distante dalla realtà).

\textbf{A Luiggi}. Sei entrato nella mia vita in punta di piedi, e piano piano ti sei ritagliato uno spazio fondamentale. Discorsi infiniti e confronti interminabili sono solo la punta dell'iceberg di quello che rappresenti ormai per me. Grazie per il tuo supporto in questo percorso e per la genuinità con cui affronti sempre ogni situazione. A Clash, comunque, rimani una pippa (\emph{skill issue}, pure a te).

\textbf{A Francesca}, il ponte che ha reso possibile tutto questo. Fermarmi a parlare con te quella mattina mi ha stravolto l'esistenza qui a Torino. Sei (in qualche modo) il collante che ha messo insieme questa banda di deficienti (sta a te valutare se questo sia un merito o meno). Grazie per ciò che fai e per metterci sempre tanto impegno.

\textbf{A Simone}. Il bene che ti voglio non lo riesco nemmeno a mettere in parole. Grazie per essere sempre di supporto a tutti, senza farlo pesare mai. Sei quella persona che c'è sempre e che mette gli altri al primo posto. La tua bontà è la tua forza più grande. Grazie per i sorrisi che ci regali. Sei quella persona di cui tutti avrebbero bisogno. Non cambiare mai.

\textbf{A Gianluca}. Ormai la tradizione degli impasti domenicali è irrinunciabile. Spero un giorno di poter essere anch'io così fortunato da farmi un sandwich con te alle tre del mattino.

\textbf{A Ilaria}, la non ingegnera del gruppo. Grazie per la spensieratezza che trasmetti e per le risate che ci regali sempre. La tua presenza rende tutto più divertente, anche se non sei un ingegnere.

\textbf{A Sara}. Anche da ultima arrivata sei riuscita a entrare a far parte di questo gruppo pazzo. Sei l'unica vera conoscenza della triennale che mi sono ritrovato anche qui a Torino, completamente a caso.

\textbf{A Jacopo}, amico di una vita. Negli anni hai saputo diventare una costante e il momento dell'anno in cui riusciamo a vederci è da sempre uno di quelli che attendo con ansia. Abbiamo già condiviso tantissimi ricordi indimenticabili e sono certo che ce ne saranno ancora moltissimi. Ti auguro il meglio per tutto quello che farai e per tutti i tuoi progetti, artistici e non.

\textbf{A Matteo}, compagno di avventure da due anni a questa parte. Dai progetti del primo anno al tirocinio in azienda, sei stato il collega e amico di cui tutti avrebbero bisogno. Ti faccio un grande in bocca al lupo per il tuo percorso. Meriti il meglio.

\textbf{A Botte Pazza} nella sua interezza. Siete diventati la mia famiglia qui a Torino. Senza il vostro sostegno tutto sarebbe stato più difficile, e questo non è altro che il risultato della vostra presenza nella mia vita. Siete semplicemente il gruppo migliore che abbia mai avuto e avervi trovato qui, lontano da casa, è la fortuna più grande che la vita mi abbia riservato finora.

Un grazie va anche al \textbf{Politecnico di Torino} e a tutti i suoi docenti, in particolare al \textbf{Prof.~Paolo Garza}, relatore di questa tesi, che ringrazio di cuore per avermi dato la possibilità di sviluppare un lavoro originale e significativo per me. Siete riusciti a trasmettermi l'amore e la passione per quello che facciamo, prima ancora di mere nozioni e concetti.

\noindent\begin{minipage}{\linewidth}\setlength{\parskip}{0.9\baselineskip}
Infine, a tutti voi, che in modi diversi avete reso possibile questo traguardo: grazie. Ogni pagina di questa tesi porta con sé un pezzo di ciascuno di voi.\par\hfill\textit{Francesco}
\end{minipage}

\par}
